% !TeX root = ../../Book1.tex
% 纲要标题：### 2.3 射影二次曲面与秩
% 纲要位置：几何学计划纲要.md，第 171 行。
% 正文状态：初稿；尚未正式审核。

% 纲要内容开始
% 
% * 齐次化及无穷远截面
% * 基域、秩和射影等价
% * 极大线性子空间及直纹族
% * 复化与实形式的不同分类信息
% 
% 纲要内容结束

\section{射影二次曲面与秩}
\label{g1:sec:0203}

仿射二次方程的三类系数在平移时互相混合，因此上一节的分类要分别讨论有无中心。齐次坐标提供了另一种看法：用新增坐标补齐次数，把所有项放进同一个二次型。这样可以直接使用合同分类，但代价是我们不再固定哪一个超平面被称作无穷远。椭圆、抛物线和双曲线的区别，恰好会在这个选择中重新出现。

\subsection{齐次化及无穷远截面}
\label{g1:subsec:0203:01}

设 \(F(x)=x^{\mathsf T}Ax+2b^{\mathsf T}x+c\) 是 \(K^n\) 上次数恰为二的多项式，其中 \(K=\mathbb R\) 或 \(\mathbb C\)，\(A\) 对称且非零。新增坐标 \(X_0\)，把原来的坐标写为 \(X=(X_1,\ldots,X_n)^{\mathsf T}\)，定义它的\term[qicihua]{齐次化}[homogenization]为
\begin{equation}
	\label{g1:eq:quadric-homogenization}
	\begin{aligned}
		F^{\mathrm h}(X_0,X)
		&=X^{\mathsf T}AX+2X_0b^{\mathsf T}X+cX_0^2\\
		&=\begin{pmatrix}X_0&X^{\mathsf T}\end{pmatrix}
		M\begin{pmatrix}X_0\\X\end{pmatrix},\\
		M&=\begin{pmatrix}c&b^{\mathsf T}\\b&A\end{pmatrix}.
	\end{aligned}
\end{equation}
它的每一项总次数都是二，所以向量整体乘以 \(t\ne0\) 时，函数值乘以 \(t^2\)；零值条件因此能在射影空间上定义。

\begin{definition}[射影二次曲面]
	\label{g1:def:projective-quadric}
	设 \(V\) 是有限维 \(K\) 向量空间，\(K=\mathbb R\) 或 \(\mathbb C\)，\(q\) 是 \(V\) 上非零二次型。由齐次方程 \(q=0\) 定义的\term[sheyingerciqumian]{射影二次曲面}[projective quadric]记为
	\[
		Q(q)=\{\,[v]\in\mathbb P(V):q(v)=0\,\}.
	\]
	讨论方程等价时，保留 \(q\) 到整体非零倍数的信息。其秩是对应对称矩阵的秩；矩阵可逆时称该二次曲面方程非退化，否则称其退化。
\end{definition}

在标准仿射图 \(X_0\ne0\) 中，用 \(x_i=X_i/X_0\)，有
\[
	F^{\mathrm h}(X_0,X)=X_0^2F(X/X_0).
\]
因此这个图中的零点恰好是原来的 \(F=0\)。在无穷远超平面 \(X_0=0\) 上，剩下的方程则是 \(X^{\mathsf T}AX=0\)。它记录二次项的零值方向，称为相对于这个仿射图的\term[wuqiongyuanjiemian]{无穷远截面}[section at infinity]。若把次数小于二的多项式也机械地补到二次，就可能多出整个无穷远超平面，例如线性方程 \(x_1=0\) 被写成 \(X_0X_1=0\)；这正是这里要求原方程次数恰为二的原因。

\begin{example}[三种实二次曲线的无穷远点]
	\label{g1:exm:conics-at-infinity}
	仿射平面中的三个方程齐次化后成为
	\[
		\begin{array}{c|c|c}
			\text{仿射方程}&\text{齐次方程}&\text{实无穷远点}\\ \hline
			x_1^2+x_2^2=1&X_1^2+X_2^2-X_0^2=0&\text{无}\\
			x_1^2-x_2^2=1&X_1^2-X_2^2-X_0^2=0&[0:1:1],\ [0:1:-1]\\
			x_1^2=x_2&X_1^2-X_0X_2=0&[0:0:1]
		\end{array}
	\]
	最后一行在无穷远超平面上的限制为平方 \(X_1^2\)，因此只给出一个实无穷远点。中间一行有两个不同的零值方向；第一行的正定二次项没有非零实零值方向。这里说的无穷远，始终相对于已经选定的 \(X_0=0\)。
\end{example}

在实数域上，还应区分齐次化与通常拓扑中的闭包。例如 \(x_1^2+1=0\) 在 \(\mathbb R^2\) 中没有实点，而齐次方程 \(X_1^2+X_0^2=0\) 在 \(\mathbb P^2(\mathbb R)\) 中有点 \([0:0:1]\)。所以齐次化的实零点集不必等于原实零点集的拓扑闭包。本节依据的是多项式和齐次坐标的计算，不把实点的逼近性质暗中放进“补入无穷远点”的说法中。

\ProseParagraphBreak

仿射坐标变换 \(x=Cy+a\) 在齐次坐标中对应可逆矩阵
\[
	T=\begin{pmatrix}1&0\\a&C\end{pmatrix}.
\]
它把 \((Y_0,Y)\) 送到 \((Y_0,aY_0+CY)\)，所以保持 \(X_0=0\)，并使 \(M\) 变为 \(T^{\mathsf T}MT\)。若允许任意可逆的 \((n+1)\) 阶矩阵，就不必再保持这个超平面，得到范围更大的射影等价关系。

\begin{proposition}[齐次化的秩]
	\label{g1:prop:homogenized-rank}
	设 \(F\) 是 \(\mathbb R^n\) 上次数恰为二的实多项式，二次项矩阵 \(A\) 的惯性为 \((p,q,r)\)，秩为 \(k=p+q\)。记齐次化矩阵为 \(M\)。
	\begin{enumerate}
		\item 若 \(F\) 有中心，中心值为 \(\kappa\)，则 \(M\) 合同于 \(\operatorname{diag}(\kappa,A)\)，其秩为 \(k\) 或 \(k+1\)，分别对应 \(\kappa=0\) 或 \(\kappa\ne0\)。
		\item 若 \(F\) 没有中心，则 \(\operatorname{rank}M=k+2\)，惯性为 \((p+1,q+1,r-1)\)。
	\end{enumerate}
\end{proposition}

\begin{proof}
	有中心时，把该中心移到原点，式\eqref{g1:eq:quadratic-translation}给出 \(F=y^{\mathsf T}Ay+\kappa\)，其齐次化矩阵正是 \(\operatorname{diag}(\kappa,A)\)。

	无中心时，\cref{g1:thm:euclidean-quadric-normal}给出 \(F=\sum_{i=1}^k\lambda_i y_i^2+2\beta z_1\)，其中 \(\beta>0\)。齐次化后，\(z_1\) 与新增坐标 \(Y_0\) 组成矩阵块
	\[
		\begin{pmatrix}0&\beta\\\beta&0\end{pmatrix},
	\]
	它在方向 \((1,1)\)、\((1,-1)\) 上分别有特征值 \(\beta\)、\(-\beta\)。所以它贡献一个正方向、一个负方向，原核中的其余 \(r-1\) 个方向仍为零。由此同时得到秩与惯性公式。
\end{proof}

这说明抛物线的二次项虽有核，齐次化的二次型却可以非退化；也说明一个有中心方程可能因中心值变成零而使齐次化的秩下降。对上一节的族 \(x^2+y^2=t\)，齐次化矩阵为 \(\operatorname{diag}(-t,1,1)\)，正好在 \(t=0\) 退化。

\subsection{基域、秩和射影等价}
\label{g1:subsec:0203:02}

对于 \(K=\mathbb R\) 或 \(\mathbb C\)，二次型 \(q\) 都对应唯一的对称双线性型
\[
	B(u,v)=\frac{q(u+v)-q(u)-q(v)}2.
\]
这与\cref{g1:def:symmetric-form}的实数情形相同。复数情形仍然对两个变量分别线性，矩阵公式是 \(u^{\mathsf T}Av\)，不含复共轭；它并不是用来测长度的复内积。根基仍定义为与所有向量配对均为零的向量组成的子空间。

\ProseParagraphBreak

在实数域上，惯性定理已经证明正负号无法通过实可逆换基消去；在复数域上，\(i^2=-1\) 却使 \(-x^2=(ix)^2\)。因此复数域上的分类应当更粗。为避免把实自伴谱定理误用到复对称矩阵上，我们直接给出所需的复数证明。

\begin{lemma}[复对称型的合同标准形]
	\label{g1:lem:complex-symmetric-normal}
	设 \(V\) 是 \(N\) 维复向量空间，\(B\) 是其上的对称双线性型，秩为 \(s\)。存在一组复基，使其矩阵为 \(\operatorname{diag}(I_s,0_{N-s})\)。两个同维复对称双线性型合同，当且仅当它们的秩相同。
\end{lemma}

\begin{proof}
	若 \(B=0\)，任意基都满足结论。若 \(B\ne0\)，则存在 \(v\) 使 \(B(v,v)\ne0\)：否则由极化公式，所有 \(B(u,v)\) 都为零。非零复数都有平方根，所以缩放 \(v\) 可得到 \(e\) 满足 \(B(e,e)=1\)。对任意 \(w\in V\)，分解
	\[
		w=B(w,e)e+\bigl(w-B(w,e)e\bigr)
	\]
	的第二项与 \(e\) 正交。并且 \(\mathbb Ce\cap e^\perp=0\)，因为 \(B(te,e)=t\)。这里 \(e^\perp\) 表示关于 \(B\) 的正交空间，故 \(V=\mathbb Ce\oplus e^\perp\)，在这个分解下矩阵为一个 \(1\) 与较低维对称型的直和。

	继续对后者归纳，直到剩下零型，就得到对角元只含 \(1\) 与 \(0\) 的矩阵。乘以可逆矩阵不改变秩，所以其中 \(1\) 的个数必为 \(s\)。这既给出数目唯一性，也说明秩相同的两个型都合同于同一个标准矩阵；反向蕴含由秩在合同变换下不变得到。
\end{proof}

\begin{theorem}[实、复射影二次方程的分类]
	\label{g1:thm:projective-quadric-classification}
	设 \(K=\mathbb R\) 或 \(\mathbb C\)，\(q,q'\) 是 \(N\) 维 \(K\) 向量空间上的非零二次型。称其射影方程等价，是指存在 \(T\in\operatorname{GL}(N,K)\) 及 \(\mu\in K\setminus\{0\}\)，使
	\[
		q'(v)=\mu q(Tv)\qquad(v\in K^N).
	\]
	在复数域上，等价的充要条件是二者秩相同。在实数域上，若惯性分别为 \((p,q_-,r)\) 与 \((p',q'_-,r')\)，等价的充要条件是
	\[
		r=r',\qquad \{p,q_-\}=\{p',q'_-\},
	\]
	其中花括号表示允许交换两个惯性指数的无序二元组。
\end{theorem}

\begin{proof}
	矩阵变化为 \(A' =\mu T^{\mathsf T}AT\)，故任何基域上秩都保持不变。复数情形由\cref{g1:lem:complex-symmetric-normal}，相同秩已足以给出合同，甚至不必再乘一个标量。

	实数情形由\cref{g1:thm:inertia}，合同保持 \((p,q_-,r)\)。乘以正数不改变这三个数；乘以负数则交换正、负惯性指数，保持根基维数。故所列条件必要。若条件成立，两组惯性相同，或把其中一个型乘以 \(-1\) 后相同；再用惯性定理的标准形便得到所需合同，故条件也充分。
\end{proof}

这里用 \(q_-\) 表示负惯性指数，是为了与二次型 \(q\) 区别；下面没有这种冲突时仍写惯性为 \((p,q,r)\)。定理分类的是方程及其基域，不只是肉眼所见的实点形状。例如在 \(\mathbb P^2(\mathbb R)\) 中，方程 \(X_0^2=0\) 与 \(X_0^2+X_1^2=0\) 的实零点集分别是一条直线与一个点；它们的秩分别为 \(1,2\)。即使两个方程的实零点集偶然相同，也不能因此认定它们按上述定义等价。

\begin{example}[椭圆、双曲线和抛物线的射影类型]
	\label{g1:exm:one-projective-conic-type}
	\cref{g1:exm:conics-at-infinity}中的三个齐次型都具有秩 \(3\)，正、负惯性指数的无序二元组均为 \(\{2,1\}\)，所以它们的实射影方程等价。第三个式子中，\(-X_0X_2\) 在 \(X_0+X_2\) 与 \(X_0-X_2\) 两个方向上给出相反符号，再加上 \(X_1^2\) 就能直接核对惯性。它们作为仿射曲线仍然不同，因为给出等价的射影变换不必保持已选的无穷远直线。区分“曲线”与“曲线加上一条指定的无穷远直线”，就能同时理解这两个结论。
\end{example}

\subsection{极大线性子空间及直纹族}
\label{g1:subsec:0203:03}

一个射影点 \([v]\) 在 \(Q(q)\) 上，只要求 \(q(v)=0\)。若希望整条射影直线、甚至更高维的射影子空间都在曲面上，就必须要求它的全部向量都满足这个条件。由\cref{g1:prop:projective-dimension}，\(k\) 维射影子空间来自 \(k+1\) 维向量子空间；因此以下问题可以完全化成双线性型问题。

\begin{definition}[全迷向子空间]
	\label{g1:def:totally-isotropic}
	设 \(V\) 是实或复向量空间，\(B\) 是对称双线性型，\(q(v)=B(v,v)\)。子空间 \(U\subseteq V\) 若满足 \(q|_U=0\)，称为\term[quanmixiangzikongjian]{全迷向子空间}[totally isotropic subspace]。等价地，\(B(u,w)=0\) 对所有 \(u,w\in U\) 成立；等价性由极化公式得到。不能真包含于更大全迷向子空间中的 \(U\) 称为极大，而可能达到的最大维数还需证明。
\end{definition}

注意，全迷向条件要求子空间内任意两个向量之间的配对也为零。两个分别迷向的向量未必张成全迷向平面：对 \(x^2-y^2\)，\((1,1)\) 与 \((1,-1)\) 都迷向，但它们的和为 \((2,0)\)，二次型值为 \(4\)。

\begin{theorem}[二次曲面中的最大线性维数]
	\label{g1:thm:quadric-linear-dimension}
	设 \(V\) 是 \(N\) 维实或复向量空间，\(q\ne0\) 是二次型，对应型为 \(B\)，根基维数为 \(r\)。
	\begin{enumerate}
		\item 若基域为实数，\(B\) 的惯性为 \((p,q_-,r)\)，则全迷向子空间的最大维数为
		\[
			r+\min(p,q_-).
		\]
		\item 若基域为复数，\(B\) 的秩为 \(s=N-r\)，则该最大维数为
		\[
			r+\lfloor s/2\rfloor.
		\]
	\end{enumerate}
	每个按包含关系极大的全迷向子空间都达到相应最大维数。若该数为 \(d>0\)，则 \(Q(q)\) 中最大射影线性子空间的维数为 \(d-1\)；若 \(d=0\)，则 \(Q(q)\) 没有基域内的点。
\end{theorem}

\begin{proof}
	先处理根基。由标准形可取直和 \(V=R\oplus W\)，其中 \(R=\operatorname{rad}B\)，\(B|_W\) 非退化。任意全迷向子空间 \(U\) 投影到 \(W\) 的像 \(\overline U\) 仍全迷向，而且
	\[
		\dim U=\dim(U\cap R)+\dim\overline U
		\le r+\dim\overline U.
	\]
	反过来，\(W\) 中的全迷向子空间与 \(R\) 直和仍全迷向。所以最大维数等于 \(r\) 加上非退化部分中的最大维数。

	在实非退化情形，标准形为 \(\lVert x\rVert^2-\lVert y\rVert^2\)，其中 \(x\in\mathbb R^p\)、\(y\in\mathbb R^{q_-}\)。对全迷向子空间，投影到任一组坐标都是单射：例如 \(x=0\) 时，零值条件给出 \(\lVert y\rVert^2=0\)，迫使 \(y=0\)。所以其维数同时不超过 \(p\) 与 \(q_-\)。取正、负标准基 \(e_i,f_j\)，向量
	\[
		e_1+f_1,\ldots,e_m+f_m,
		\qquad m=\min(p,q_-),
	\]
	彼此配对均为零，且线性无关，故上界可达到。

	在复非退化的 \(s\) 维空间中，若 \(U\) 全迷向，则 \(U\subseteq U^\perp\)。这里 \(\dim U^\perp=s-\dim U\)：非退化性使 \(w\mapsto B(w,-)\) 成为到对偶空间的同构，而任意 \(U\) 上的线性函数都可以先扩充一组基，再延拓到全空间；因此 \(w\mapsto B(w,-)|_U\) 满射到 \(U^*\)，其核就是 \(U^\perp\)。秩与零化度公式给出所需维数。于是 \(2\dim U\le s\)。在\cref{g1:lem:complex-symmetric-normal}的标准基中，
	\[
		e_1+ie_2,\ e_3+ie_4,\ldots,
		e_{2m-1}+ie_{2m},\qquad m=\lfloor s/2\rfloor,
	\]
	构成线性无关的全迷向向量组，故复数情形的上界也可达到。

	最后证明极大子空间不可能停在更小的维数。极大全迷向子空间必包含 \(R\)，否则加上根基中的一个新方向仍然全迷向；此时它等于 \(R\oplus(U\cap W)\)。于是只需在非退化部分中证明结论，仍把其中的全迷向子空间记为 \(U\)。若 \(u\in U\) 非零，非退化性给出 \(v\) 满足 \(B(u,v)=1\)。以
	\[
		v'=v-\frac{B(v,v)}2u
	\]
	替代 \(v\)，就有 \(B(v',v')=0\)、\(B(u,v')=1\)。平面 \(H=\operatorname{span}(u,v')\) 的配对矩阵为 \(\left(\begin{smallmatrix}0&1\\1&0\end{smallmatrix}\right)\)，因而非退化。任意 \(w\) 都可按
	\[
		w=B(w,v')u+B(w,u)v'
		+\bigl(w-B(w,v')u-B(w,u)v'\bigr)
	\]
	分成 \(H\) 与 \(H^\perp\) 两部分，且交集为零。\(H^\perp\) 上的限制仍非退化：若其中一个向量与整个 \(H^\perp\) 正交，它也与 \(H\) 正交，便与整个空间正交，故只能为零。若 \(w\in U\)，则 \(B(w,u)=0\)，所以 \(U=Ku\oplus(U\cap H^\perp)\)。

	重复这一步，若 \(\dim U=k\)，便剥离出 \(k\) 个这样的二维平面，余下非退化空间 \(W_0\) 与它们正交。在实数域中每剥离一次，正、负惯性指数各减一，所以若 \(k<\min(p,q_-)\)，\(W_0\) 仍有一正一负两个方向；它们可组成新的迷向向量。在复数域中，若 \(k<\lfloor s/2\rfloor\)，则 \(\dim W_0\ge2\)，也可由标准形的前两个方向组成 \(e_1+ie_2\)。这些新向量都与原 \(U\) 正交，使 \(U\) 真正扩大，违背极大性。这也覆盖 \(U=0\) 而上界为正的情形。

	最后，非零全迷向子空间 \(U\) 与曲面内的射影子空间 \(\mathbb P(U)\) 恰好对应，其维数减少一。若不存在非零全迷向子空间，连一个迷向向量都不存在，曲面便没有相应基域的射影点。
\end{proof}

定理中的两个维数计数揭示了实、复几何的差别。例如实球面齐次化后的型为 \(X_1^2+X_2^2+X_3^2-X_0^2\)，惯性为 \((3,1,0)\)。全迷向子空间最大只有一维，所以球面射影化后没有实射影直线。允许复坐标以后，同一个矩阵秩为 \(4\)，最大全迷向维数变为二，因而出现复射影直线。

\ProseParagraphBreak

对于惯性 \((2,2,0)\) 的实型，最大射影线性子空间就是直线。下面的模型不仅表明直线存在，而且把它们全部写出。这个表示也将在\secref{g1:sec:0206}中成为一个射影映射的例子。

\begin{proposition}[二次曲面的两个直纹族]
	\label{g1:prop:quadric-two-rulings}
	设 \(K=\mathbb R\) 或 \(\mathbb C\)，在 \(\mathbb P^3(K)\) 中考虑
	\[
		Q:\ X_0X_3-X_1X_2=0.
	\]
	每个点唯一写为
	\begin{equation}
		\label{g1:eq:quadric-rank-one-param}
		[X_0:X_1:X_2:X_3]
		=[su:sv:tu:tv],
		\qquad [s:t],[u:v]\in\mathbb P^1(K).
	\end{equation}
	固定 \([s:t]\) 而让 \([u:v]\) 变化，或固定 \([u:v]\) 而让 \([s:t]\) 变化，各得到一个直纹族。这两族包含 \(Q\) 上的全部射影直线；同族不同直线互不相交，异族任意两条直线恰交于一点。过 \(Q\) 上的每一点恰有两条直线，每族各一条。
\end{proposition}

\begin{proof}
	把四个坐标排成矩阵
	\[
		Z=\begin{pmatrix}X_0&X_1\\X_2&X_3\end{pmatrix}.
	\]
	方程是 \(\det Z=0\)，而射影坐标要求 \(Z\ne0\)，所以 \(Z\) 的秩恰为一。选一个非零列，其他列都是它的倍数，便可写成 \(Z=a b^{\mathsf T}\)，其中非零列向量 \(a=(s,t)^{\mathsf T}\)、\(b=(u,v)^{\mathsf T}\)。它的列空间决定 \([a]\)，行空间决定 \([b]\)，故这两个射影点都唯一。分别改变 \(a,b\) 的代表，只把 \(Z\) 整体乘以非零数，所以式\eqref{g1:eq:quadric-rank-one-param}良定义。

	固定非零 \(a\) 后，映射 \(b\mapsto ab^{\mathsf T}\) 线性且单射，因此其像是二维向量子空间，射影化后正是一条直线；固定 \(b\) 同理。分解的唯一性说明同族不同直线没有公共点，而分别固定 \([a]\)、\([b]\) 的异族两条线只在 \([ab^{\mathsf T}]\) 相交。

	还须排除第三种直线。设 \(L\subseteq Q\) 为任意射影直线，取它对应二维向量空间的一组基 \(ab^{\mathsf T}\)、\(cd^{\mathsf T}\)。对所有 \(\tau\in K\)，矩阵 \(ab^{\mathsf T}+\tau cd^{\mathsf T}\) 的行列式为零。按两列展开，得到
	\[
		\det\bigl(ab^{\mathsf T}+\tau cd^{\mathsf T}\bigr)
		=\tau\det(a,c)\det(b,d).
	\]
	取 \(\tau=1\)，可知 \(a,c\) 成比例或 \(b,d\) 成比例；二者不能同时成比例，否则两基向量相关。在第一种情形，整条 \(L\) 属于固定 \([a]\) 的族；在第二种情形，属于固定 \([b]\) 的族。全部断言随之成立。
\end{proof}

在图 \(X_0\ne0\) 中令 \(x=X_1/X_0\)、\(y=X_2/X_0\)、\(z=X_3/X_0\)，方程成为 \(z=xy\)，两个直纹族正是上一节画出的固定 \(x\) 与固定 \(y\) 的直线族。不过射影曲面还包含无穷远处的两条线：把 \(X_0=0\) 代入，得到 \(X_1X_2=0\)，即 \(X_0=X_1=0\) 与 \(X_0=X_2=0\)。它们也各属于一个直纹族，交于 \([0:0:0:1]\)，因此射影描述补全了仿射图没有显示的直线。

\subsection{复化与实形式的不同分类信息}
\label{g1:subsec:0203:04}

一个实二次方程可以保持系数不变而允许复数坐标，这称为把方程\term[fuhua]{复化}[complexification]。矩阵的秩不变，因为某阶子式是否非零不因把实数看成复数而改变；惯性的正负号却不再是复换基的不变量。乘以 \(i\) 就能把一个正平方改成负平方，因而所有相同秩的复二次型归入同一类。

\ProseParagraphBreak

为说明丢掉了哪些信息，\cref{g1:tab:real-quadric-forms}比较 \(\mathbb P^3(\mathbb R)\) 中的三种满秩方程。它们的复化都具有秩 \(4\)，所以由\cref{g1:thm:projective-quadric-classification}彼此复射影等价。

\begin{table}[htbp]
	\centering
	\caption{同一个复二次曲面类型的不同实形式}
	\label{g1:tab:real-quadric-forms}
	\begin{tabularx}{\linewidth}{@{}l c X@{}}
		\toprule
		实齐次方程 & 惯性 & 实几何 \\
		\midrule
		\(X_0^2+X_1^2+X_2^2+X_3^2=0\) & \((4,0,0)\) & 没有实点 \\
		\(X_1^2+X_2^2+X_3^2-X_0^2=0\) & \((3,1,0)\) & 有实点，没有实直线 \\
		\(X_0X_3-X_1X_2=0\) & \((2,2,0)\) & 有两个实直纹族 \\
		\bottomrule
	\end{tabularx}
\end{table}

第一行中，正定性排除了非零实零值向量；第二行在 \(X_0=1\) 中就是单位球面，维数定理已经说明它不含实直线；第三行则由两个符号相反的乘积各贡献一个正、一个负方向，且直纹由\cref{g1:prop:quadric-two-rulings}完全描述。这三种实方程可称为该复射影类型的不同\term[shixingshi]{实形式}[real forms]：它们的复化等价，但实数域上的方程并不等价。这里的名称只用于这项具体比较，不预设一般实形式理论。

\ProseParagraphBreak

由实到复并非简单地给图上每一点再加一个虚部。实点的存在、实直线的存在和实惯性都可能在复化后不再可辨；保存下来的秩则仍然控制根基维数及最大复线性子空间维数。后面讨论极性时，非退化性还会决定一个点能否唯一对应一个超平面。因此，在说“二次曲面由秩分类”以前，必须先说明基域；在说“这个曲面有直纹”以前，也必须说明这些直线在哪个基域上定义。
